% NeuroPower Protocol — arXiv / IEEE draft skeleton (Task 5)
\documentclass[conference]{IEEEtran}
\usepackage{cite}
\usepackage{amsmath,amssymb,amsfonts}
\usepackage{graphicx}
\usepackage{url}
\usepackage{booktabs}

\begin{document}

\title{NeuroPower: An Open Layered Protocol for Event-Driven Neuromorphic Memory--Inference Co-Design on Wearables}

\author{
\IEEEauthorblockN{NeuroPower Protocol Contributors}
\IEEEauthorblockA{Open collaboration --- see project repository}
}

\maketitle

\begin{abstract}
Proprietary neuromorphic APIs fragment the wearable AI ecosystem and make sub-watt continuous inference difficult to port. We present \emph{NeuroPower}, an open layered protocol (Physical through Application) for spike-based memory and inference, with a synthesizable RISC-V + LIF accelerator reference, offline INT4 training export, and event-driven power models targeting $<$1\,W peak, $<$100\,mW idle, and $\geq$12\,h on a 200\,mAh cell.
\end{abstract}

\begin{IEEEkeywords}
neuromorphic computing, open hardware, wearable AI, RISC-V, spiking neural networks
\end{IEEEkeywords}

\section{Introduction}
% Problem: Loihi / Akida / TrueNorth lock-in; wearable battery limits.
Wearable always-on AI is constrained by energy and by incompatible neuromorphic interfaces. NeuroPower defines an interoperable stack and a free reference implementation verified in simulation (no tape-out required for protocol validation).

\section{Protocol}
\subsection{Physical and Link}
SPI Mode~0 plus asynchronous \texttt{SPIKE\_IN}/\texttt{SPIKE\_OUT}, with a 5-byte Link packet and CRC-8~\cite{neuropower-spec}.

\subsection{Network and Application}
Up to 16 cores $\times$ 256 neurons; C API \texttt{np\_init}/\texttt{np\_load\_weights}/\texttt{np\_inference}.

\section{Reference Implementation}
Ibex/VexRiscv-class RV32I host bus, 256-LIF array with INT4 synapses, \texttt{.npw} weight streaming, and clock gating via \texttt{power\_ctrl}.

\section{Evaluation}
Offline keyword-spotting stand-in exceeds 90\% INT4 membrane accuracy on synthetic spikes; power scripts estimate $\geq$12\,h autonomy at 10\% inference duty on 200\,mAh.

\section{Conclusion}
NeuroPower offers a practical open baseline for neuromorphic wearable co-design. Community review and IEEE standardization pathways (e.g., related working groups) are future work.

\begin{thebibliography}{1}
\bibitem{neuropower-spec} NeuroPower Protocol, ``Specification v0.1,'' project repository, 2026.
\end{thebibliography}

\end{document}
